\section{Matched and stratified contrasts}
\label{sec:contrasts}

The marginal means of Section~\ref{sec:marginals} average over an unbalanced grid. This
section reports the same axes as matched within-design contrasts, formed only from cell
pairs that agree on every other factor and on seed, and then as fully stratified contrasts
recomputed inside every combination of benchmark, capacity and reasoning rung.

\subsection{Matched adjacent contrasts}

Table~\ref{tab:contrasts-accuracy} gives all 92 matched accuracy contrasts, and
Figure~\ref{fig:contrast-forest} displays them on one page.
Appendix~\ref{app:marginals} carries the same contrasts for the ten other metrics.

\intable{tab_contrasts_accuracy}

\begin{figure}[H]
\centering
\includegraphics[width=\linewidth,height=0.66\textheight,keepaspectratio]{fig_contrast_forest.pdf}
\caption{All matched adjacent-level accuracy contrasts, grouped by axis and faceted by
benchmark. Bars are hierarchical bootstrap percentile intervals.}
\label{fig:contrast-forest}
\end{figure}

The capacity ladder dominates. Adjacent steps are $+0.043$, $+0.153$ and $+0.181$ from
0.6B to 1.7B on GPQA, MMLU-Pro and TruthfulQA respectively, and the endpoint contrast from
0.6B to 14B is $+0.274$ $[0.253, 0.294]$, $+0.428$ $[0.418, 0.439]$, $+0.437$ $[0.421,
0.452]$ and $+0.322$ $[0.299, 0.345]$ across the four benchmarks. The 32B against 14B
contrast is $+0.002$ $[-0.011, 0.015]$, $+0.011$ $[0.006, 0.018]$, $+0.006$ $[-0.001,
0.012]$ and $-0.016$ $[-0.031, -0.002]$: two of four intervals contain zero and the
mathematical one is negative.

The reasoning ladder shows the same saturation. Against thinking disabled, b2048 gives
$+0.087$ $[0.070, 0.103]$, $+0.135$ $[0.120, 0.149]$, $+0.092$ $[0.071, 0.115]$ and
$+0.053$ $[0.037, 0.071]$; b8192 gives $+0.098$, $+0.148$, $+0.087$ and $+0.097$; and
unlimited gives $+0.093$, $+0.150$, $+0.090$ and $+0.092$. The unlimited rung is
indistinguishable from b8192 on three benchmarks and slightly worse on two.

Topology contrasts are an order of magnitude smaller than capacity contrasts.
Independent against single agent gives $+0.010$ $[0.004, 0.018]$, $+0.024$ $[0.018,
0.030]$, $+0.011$ $[0.006, 0.016]$ and $+0.076$ $[0.062, 0.091]$. Decentralized against
independent gives $+0.007$ $[0.000, 0.013]$, $+0.016$ $[0.011, 0.021]$, $+0.010$ $[0.006,
0.015]$ and $+0.009$ $[-0.000, 0.017]$. Centralized against independent gives $+0.002$
$[-0.005, 0.009]$, $+0.008$ $[0.001, 0.016]$, $+0.003$ $[-0.004, 0.009]$ and $-0.071$
$[-0.090, -0.054]$. Centralized against decentralized gives $-0.004$, $-0.011$ $[-0.018,
-0.005]$, $-0.003$ and $-0.086$ $[-0.105, -0.067]$.

The mathematical benchmark is where the topologies separate most sharply. Independent
voting gains $+0.076$ there, the largest topology effect anywhere in the sweep, while
centralized orchestration loses $0.071$ against the same baseline. A single orchestrator
synthesising two sub-agent solutions is substantially worse on free-form mathematics than
simply taking the majority of three independent attempts.

Table~\ref{tab:orchestration} in Appendix~\ref{app:fits} repeats the topology contrasts
across all six metrics for which they are computed, including coordination efficiency,
error amplification and the cost proxy.

Moving from independent to decentralized costs 0.067, 0.104, 0.104 and 0.112 of
coordination efficiency on the four benchmarks while buying between 0.007 and 0.016 of
accuracy. Debate raises vote calibration error against independent voting by 0.023, 0.007,
0.019 and 0.016.

\subsection{Stratified contrasts}

A pooled contrast can conceal a sign reversal by regime. Every non-capacity, non-reasoning
contrast was therefore recomputed inside each benchmark, capacity and reasoning stratum.
Table~\ref{tab:stratum-digest} summarises the result, and
Figure~\ref{fig:stratum-caterpillar} shows every individual stratum estimate.

\intable{tab_stratum_digest}

\begin{figure}[H]
\centering
\includegraphics[width=\linewidth,height=0.48\textheight,keepaspectratio]{fig_stratum_caterpillar.pdf}
\caption{Every stratified accuracy contrast, sorted within contrast, coloured by
benchmark, with the pooled matched estimate as a horizontal line. The width of each cloud
around the pooled line is the extent to which the pooled number is a regime average rather
than a general effect.}
\label{fig:stratum-caterpillar}
\end{figure}

The digest reads as follows. The three-agent against one-agent accuracy contrast is the
most robust: 81 percent of its 117 strata have intervals excluding zero and 93 percent
share the pooled sign. Prompt complexity is next at 64 and 69 percent over 119 strata.
Context sharing is the least stable of the four, with 48 percent of strata excluding zero
and 73 percent sharing the pooled sign, and an interquartile range spanning zero.

Two structural cautions apply to the stratified tables. The minimum matched-pair count in
several strata is one, so individual cells of the stratified tables are localisation
devices rather than inferential statements. And the topology contrasts in particular become
thin after splitting by capacity and reasoning, since topology, capacity and reasoning are
crossed but not deeply replicated.

\subsection{Context sharing by regime}

Context sharing is the axis for which the stratified view most changes the conclusion.
Averaged within benchmark, the controlled plus-CoT accuracy effect is $+0.018$ on GPQA
with 24 of 29 strata positive, $+0.016$ on MMLU-Pro with 20 of 30 positive, $+0.008$ on
TruthfulQA with 15 of 30 positive, and $+0.050$ on MATH with 28 of 30 positive.

The banded view identifies where the effect lives:

\begin{center}
\small
\begin{tabular}{llrl}
\toprule
Capacity band & Reasoning band & Mean effect & Positive strata \\
\midrule
0.6B, 1.7B & off, b512 & $+0.031$ & 12 of 16 \\
0.6B, 1.7B & b2048 and above & $+0.024$ & 17 of 24 \\
4B, 8B & off, b512 & $+0.045$ & 16 of 16 \\
4B, 8B & b2048 and above & $+0.011$ & 18 of 24 \\
14B, 32B & off, b512 & $+0.032$ & 13 of 16 \\
14B, 32B & b2048 and above & $+0.006$ & 11 of 23 \\
\bottomrule
\end{tabular}
\end{center}

The mid-capacity low-reasoning band is unanimous: all 16 strata positive with a mean effect
of $+0.045$. The large-capacity high-reasoning band is the reverse, with 11 positive and 12
negative strata and a mean of $+0.006$. Peer rationales help when the receiving model has
capacity to use them but has not been given enough thinking budget to generate equivalent
content itself. Once the model has both capacity and budget, the peer's reasoning adds
nothing and the effect becomes noise around zero.

The calibration cost does not share this regime dependence. Controlled vote calibration
error rises with plus-CoT in 94 of 119 strata with a mean of $+0.022$, and the vote delta
rises in 71 of 89 strata with a mean of $+0.024$. Exposing rationales makes agreement a
worse confidence signal almost everywhere, including in the strata where it does not
improve accuracy.

Appendix~\ref{app:stratified} contains the complete stratified tables for accuracy, vote
calibration error and the cost proxy on all four stratified axes, 2{,}094 estimates in
total.

\subsection{Other controlled axes}

Prompt complexity remains positive overall under stratification but is less uniform than
the marginal picture suggests, with GPQA and MATH weakest after fixing capacity and
reasoning, and the large-capacity high-reasoning band shrinking to $+0.011$ with more
negative than positive strata.

Decentralized against independent is usually positive and small. Almost every topology
stratum rests on fewer than three matched design pairs after the capacity and reasoning
split, so these should be read as localisation only.

Centralized against decentralized remains negative and is worst on the mathematical
benchmark, where the controlled mean is $-0.087$ with 27 negative and zero positive
strata. This is the most one-sided stratified result in the sweep.

The three-agent family contrast is positive in 109 of 117 controlled strata, which is a
useful consistency check on the pooled estimate but does not substitute for varying agent
count directly.
